\chapter{Related work}

\section{Code clones}

Code clones are similar or identical code fragments that appear in different parts of a codebase.
They can be classified into different types based on their similarity \cite{roy09}.

\begin{codepair}[Type 1 clones]
\label{prg:type1clones}
\begin{lstlisting}[language=C++]
int add3(int a) {
    // sum of parameter and 3
    return a + 3;
}
\end{lstlisting}
\codepairsep
\begin{lstlisting}[language=C++]
int add3(int a) {
  return a + 3;
}
\end{lstlisting}
\end{codepair}

Type 1 clones are identical code fragments that differ only in whitespace and comments, see Program \ref{prg:type1clones}.

\begin{codepair}[Type 2 clones]
\label{prg:type2clones}
\begin{lstlisting}[language=C++]
int add3(int a) {
    // sum of parameter and 3
    return a + 3;
}
\end{lstlisting}
\codepairsep
\begin{lstlisting}[language=C++]
float add4(float x) {
  return x + 4.0f;
}
\end{lstlisting}
\end{codepair}

Type 2 clones, apart from whitespace and comments, also allow for renaming of identifiers and change in literals and types, see Program \ref{prg:type2clones}.

\begin{codepair}[Type 3 clones]
\label{prg:type3clones}
\begin{lstlisting}[language=C++]
int add3(int a) {
    // sum of parameter and 3
    return a + 3;
}
\end{lstlisting}
\codepairsep
\begin{lstlisting}[language=C++]
float add4(float x) {
  auto result = x + 4.0f;
  return result;
}
\end{lstlisting}
\end{codepair}

Type 3 clones might have some statements added, modified, or removed, but still show a significant similarity in their structure and functionality, see Program \ref{prg:type3clones}.

Type 4 clones are code fragments that perform the same functionality but are implemented differently.
An example of type 4 clones would be two different implementations of a sorting algorithm.



\section{Tree edit distance}

A tree is a data structure that consists of nodes connected by edges, where each node has a label and zero or more child nodes.
Tree edit distance is a measure of similarity between two trees, $T_1$ and $T_2$.
It is defined as the minimum cost of a sequence of edit operations that transforms $T_1$ into $T_2$.
The edit operations typically include insertion, deletion, and relabeling.
The tree edit distance has applications in various fields, including bioinformatics.
